\documentclass[12pt]{article}
\usepackage[T1]{fontenc}
\usepackage[utf8]{inputenc}
\usepackage[a4paper,margin=1.7cm]{geometry}
\usepackage{booktabs}
\usepackage{array}

\begin{document}

\begin{table}[htbp]
    \centering
    \caption{Workload profile shared by adaptation test rounds 1--3}
    \label{tab:adaptation-workload}
    \renewcommand{\arraystretch}{1.25}
    \begin{tabular}{p{3.1cm} p{2.6cm} c p{7.2cm}}
        \toprule
        Stage & Elapsed time & Target users & Observed ramp in the CSVs \\
        \midrule
        Initial load & 0--30 s & 10 & Reached 10 users at 4--5 s. \\
        Medium load & 30--120 s & 20 & Reached 20 users at 39--40 s. \\
        Peak load & 120--600 s & 40 & Reached 40 users at 139--140 s. \\
        Scale-down & 600--660 s & 20 & Reduction began at 601--602 s and reached 20 users at 604--605 s. \\
        \bottomrule
    \end{tabular}

    \vspace{0.5em}
    \begin{minipage}{0.92\linewidth}
        \footnotesize
        \textit{Note.} The target user-count profile is the same in rounds 1, 2, and 3, so it is shown once. The small ramp-time differences are based on the recorded \texttt{User Count} values in each round's \texttt{stats\_history.csv}.
    \end{minipage}
\end{table}

\end{document}
